\chapter{Protocole experimental}

\section{Jeux de donnees et splits}
Le notebook final utilise:
\begin{itemize}
  \item train: \textbf{3193} exemples,
  \item validation: \textbf{798} exemples,
  \item test: \textbf{908} exemples.
\end{itemize}
Le split est derive du train officiel avec \texttt{VAL\_RATIO = 0.2} et \texttt{seed = 42}, puis le test officiel est conserve pour l'evaluation finale.

\section{Configuration du notebook final}
Configuration extraite de \texttt{Code/pointnet\_notebook.ipynb}:
\begin{itemize}
  \item modele: \texttt{PointNetFull},
  \item augmentations train: \texttt{[normalize, rot\_z, noise, shuffle, occlusion]},
  \item augmentation test: \texttt{[normalize]},
  \item \texttt{MAX\_EPOCHS = 150},
  \item early stopping active (monitor \texttt{val\_acc}, patience \texttt{50}, min\_delta \texttt{0.0}),
  \item scheduler de taux d'apprentissage actif.
\end{itemize}

\section{Campagnes automatisees complementaires}
En parallele du notebook, des campagnes de benchmark ont ete lancees avec le script robuste.
Table~\ref{tab:campaigns} resume les runs principaux exploites dans ce rapport.

\begin{table}[H]
\centering
\caption{Campagnes experimentales utilisees.}
\label{tab:campaigns}
\begin{tabular}{p{0.22\linewidth}p{0.15\linewidth}p{0.57\linewidth}}
\toprule
Campagne & Date & Usage dans le rapport \\
\midrule
C1 (20260208\_145104) & 08/02/2026 & benchmark de 10 augmentations (seed 42, 25 epochs) \\
C2 (20260211\_081948) & 11/02/2026 & benchmark robuste (4 experiences, seed 1, max 50 epochs, early stop) \\
C3 (20260211\_143359) & 11/02/2026 & run cible \texttt{plus\_occlusion, seed\_0} (max 150 epochs, monitor val\_loss) \\
\bottomrule
\end{tabular}
\end{table}

\section{Metriques reportees}
Selon les sorties disponibles, on suit principalement:
\begin{itemize}
  \item \texttt{best\_val\_acc} et \texttt{best\_val\_loss},
  \item \texttt{best\_eval\_acc} quand presente,
  \item \texttt{final\_test\_acc} pour le notebook final,
  \item matrices de confusion et courbes d'apprentissage.
\end{itemize}
Cette distinction est importante pour eviter de melanger validation interne et evaluation finale.
